\documentclass{article}

\begin{document}
    \begin{center}
        \Large \textbf{Challenge 4: Module und Environments} \\ [6pt]
        \normalsize
        Autor: Maximilian Jakob Maag B.Sc. \\
        Version: 1.0
    \end{center}
    \section{Intro}
    Wie Sie bereits bemerkt haben, hat sich viel Terraform-Code wiederholt, während Sie die letzten Challenges absolviert haben.
    In einer produktiven Umgebung wäre dieses Vorgehen zu zeitaufwändig und fehleranfällig. \\ \\
    Terraform bietet die Möglichkeit, sogenannte Module zu definieren, z. B. ein Modul für einen generischen Docker‑Stack mit einer Firewall und einem Volume.
    Dieses Modul kann anschließend aufgerufen werden. Wird das Modul abgeändert, werden alle Instanzen des Moduls ebenfalls geändert. \\
    Für ein Modul müssen folgende Dinge definiert werden: \\
    \begin{itemize}
        \item Die Variablen für das Modul
        \item Der Output des Moduls
        \item Der Terraformcode des Moduls  selbst
    \end{itemize}
    Das Modul kann anschließend in einer Environment aufgerufen werden, um deployed zu werden. \\
    \section{Aufgaben}
    In dieser Übung lernen Sie Module kennen. Es ist ihr Ziel wiederverwendbaren Code zu schreiben. Außerdem sollen Sie für unterschiedliche Zwecke unterschiedliche Umgebungen deployen.
    Gängig sind eine Umgebung für Production, eine Testumgebung Staging, eine Entwicklungsumgebung ohne Regeln zum Probieren und Experimentieren (Dev). \\ \\
    Spätestens ab dieser Challenge wird der Einsatz von Git dringend empfohlen. Sie werden viele unterschiedliche Dateien handhaben in mehreren Versionen.
    Damit Sie den Überblick nicht verlieren und ein Rollback auf eine funktionierende Version durchführen können, legen Sie für diese Challenge am besten ein Git‑Repository an.
    Ein lokales Repository ist völlig ausreichend.
    \subsection{Aufgabe 1}
        Lesen Sie die Terraform-Dokumentation zu Modules. Erarbeiten Sie auch mithilfe einer KI eine Struktur für Ihren Terraform-Code.
        Legen Sie fest, welche Directories und Subdirectories Sie verwenden möchten. Im Idealfall organisieren Sie diese in einem Repository.
        Bedenken Sie, dass Sie sowohl Module als auch Umgebungen anlegen müssen. Ein Modul kann von mehreren Umgebungen verwendet werden. \\ \\
        Richten Sie anschließend mindestens eine Umgebung für produktive Anwendungen, Eine Testumgebung und eine Entwicklungsumgebung ein.
        Prod und Staging sollen in der Cloud deployed werden. Dev sollte allerdings aus Kostengründen nur lokal auf Ihrem Computer deployed werden. \\ \\
        Deployen Sie auch alle anderen Umgebungen aus Ihrem erarbeiteten Konzept. \\ \\
        Deployen Sie je Umgebung mindestens drei Anwendungen. Beispiele wären Mailcow, Nextcloud, Owncloud, Opencloud, Portainer, Mailu, Matrix, ein eigener Git‑Server.

    \subsection{Aufgabe 2}
        Lesen Sie Ihren Terraformcode genau. Suchen Sie nach wiederholten Mustern. Entscheiden Sie in welche Module Sie Ihren code unterteilen wollen. Welche Anwendungen kommen in allen Umgebungen vor?
        Was brauchen Sie immer wieder? \\ \\
        Erarbeiten Sie mindestens 3 geeignete Module und rufen Sie diese in allen Umgebungen mindestens einmal auf. \\ \\
        Falls Sie Git verwenden: Ermitteln Sie Ihre Lines of Code für und nach dieser Änderung. Vergleichen Sie diese und erläutern Sie die Bedeutung dieser Zahlen.
        Was bedeutet das für Ihren Arbeitsalltag?

    \subsection{Aufgabe 3}
        Ändern Sie nun etwas an Ihrem Modul. Zum Beispiel einen Default‑Wert. Falls Sie VPS deployen, fügen Sie auf Modulebene eine Firewall hinzu.
        Deployen Sie Ihre Staging-Umgebung. Vergleichen Sie den Zustand dieser Umgebung mit Prod und Dev. Warum hat sich der Zustand dieser Umgebungen nicht geändert?
        Was würde sich verändern, wenn Sie prod ebenfalls deployen? \\ \\
        Führen Sie nach dem Ende der Challenge terraform destroy in allen Umgebungen aus.
\end{document}